
\documentclass[12pt,a4paper]{article}

\usepackage[utf8]{inputenc}
\usepackage[T1]{fontenc}
\usepackage[spanish]{babel}
\usepackage{geometry}
\usepackage{graphicx}
\usepackage{longtable}
\usepackage{array}

\geometry{margin=2.5cm}

\title{TA-IND-01: IEEE 29148 y tipos de requisitos de un sistema conocido}
\author{Mayra Lucila Córdova Carreño}
\date{Noviembre de 2026}

\begin{document}

\maketitle

\section{Introducción}

La ingeniería de requisitos permite identificar,
analizar y documentar las necesidades que debe
satisfacer un sistema de software. Para esta
actividad se seleccionó Wattpad, una plataforma
digital de lectura y publicación de historias.

El análisis se desarrollará tomando como
referencia la norma ISO/IEC/IEEE 29148:2018,
que establece orientaciones para la ingeniería
de requisitos durante el ciclo de vida de
sistemas y software.

\section{Descripción del sistema seleccionado}

Wattpad es una plataforma digital orientada
a la lectura, escritura y publicación de
historias. Su propósito es permitir que los
usuarios accedan a obras literarias de
diferentes géneros y que los escritores
compartan sus propias creaciones.

Entre sus funcionalidades se encuentran la
búsqueda de historias por título o autor,
la lectura organizada por capítulos,
el almacenamiento de historias en una
biblioteca personal y la interacción mediante
votos y comentarios.

También permite crear, publicar y editar
historias, recibir notificaciones y modificar
la información del perfil.

El sistema está dirigido principalmente a
lectores y escritores que acceden a la
plataforma desde dispositivos móviles o
navegadores web.

Para esta actividad se analizarán las
funciones de exploración, lectura,
interacción, publicación y gestión de cuenta,
tomando como referencia la experiencia
de uso del sistema.


\section{Resumen analitico de ISO/IEC/IEEE 29148:2018}

\subsection{Alcance y fundamentos}

La norma ISO/IEC/IEEE 29148:2018 establece
orientaciones para los procesos de ingenieria
de requisitos durante el ciclo de vida de
sistemas y software. Su alcance incluye la
definicion, documentacion y gestion de
requisitos, asi como la estructura de los
productos de informacion relacionados.

En su clausula 3, la norma define conceptos
que permiten utilizar una terminologia
consistente. Entre ellos se encuentran
requisito, parte interesada, restriccion,
verificacion, validacion y trazabilidad.
Estas definiciones son fundamentales para
distinguir las necesidades de los usuarios
de las caracteristicas que debera satisfacer
un sistema.

\subsection{Requisitos bien formulados}

La clausula 5.2.4 explica que un requisito
bien formulado debe expresar una necesidad
de manera comprensible. Su redaccion debe
identificar el sujeto, establecer la
capacidad o restriccion requerida e
incorporar las condiciones necesarias para
comprobar su cumplimiento.

En una especificacion textual, la norma
recomienda utilizar expresiones obligatorias
claramente diferenciadas de las
recomendaciones y descripciones. Tambien
destaca la importancia de evitar decisiones
de diseno innecesarias.

En Wattpad, por ejemplo, es mas adecuado
establecer que el sistema debera mostrar
los resultados de una busqueda por titulo
que imponer una tecnologia especifica
para implementar esa busqueda.

\subsection{Caracteristicas individuales}

La clausula 5.2.5 establece nueve
caracteristicas para los requisitos
individuales: necesario, apropiado,
no ambiguo, completo, singular,
factible, verificable, correcto
y conforme.

Estas caracteristicas permiten evaluar
la calidad de cada enunciado. Un requisito
singular expresa una sola necesidad,
mientras que uno verificable permite
demostrar objetivamente su cumplimiento.

Por ejemplo, la expresion ``el sistema
debera ser rapido'' es ambigua y dificil
de verificar. En cambio, especificar un
tiempo maximo de respuesta bajo condiciones
de prueba definidas permite establecer
un criterio de aceptacion.

\subsection{Caracteristicas del conjunto}

La clausula 5.2.6 identifica cinco
caracteristicas del conjunto de requisitos:
completo, consistente, factible,
comprensible y susceptible de validacion.

Su aplicacion implica revisar que los
requisitos no se contradigan, utilicen
terminos uniformes y representen las
necesidades de las partes interesadas.
En el analisis de Wattpad, esto exige
mantener coherencia entre las funciones
de busqueda, lectura, biblioteca y
participacion de los usuarios.

\subsection{Lenguaje de los requisitos}

La clausula 5.2.7 recomienda expresar
que necesita realizar el sistema sin
imponer innecesariamente como debe
implementarse. Tambien advierte sobre
terminos subjetivos, comparaciones vagas,
pronombres imprecisos y expresiones
absolutas que dificultan la verificacion.

Por esta razon, los requisitos de Wattpad
deben formularse mediante acciones
observables y condiciones medibles,
evitando expresiones como ``facil'',
``muy rapido'' o ``siempre disponible''
cuando no se define su significado.

\subsection{Atributos y clasificacion}

La clausula 5.2.8 presenta atributos
que facilitan la gestion de requisitos,
como identificacion, version, responsable,
prioridad, riesgo, justificacion,
dificultad y tipo.

La clasificacion por tipos permite
distinguir requisitos funcionales,
de interfaz, de proceso, de calidad,
de usabilidad y de factores humanos.

Para Wattpad, los requisitos funcionales
describen capacidades como leer capitulos
o registrar votos. Los requisitos de
calidad, por su parte, permiten proponer
objetivos medibles de rendimiento,
fiabilidad y seguridad.

En conjunto, la norma proporciona
criterios para transformar necesidades
en requisitos claros y evaluables,
reduciendo interpretaciones incorrectas
durante el desarrollo del software.


\section{Glosario de terminos}

Los siguientes conceptos se seleccionaron
de la clausula 3.1 de la norma
ISO/IEC/IEEE 29148:2018.
Las definiciones se presentan mediante
parafrasis y se acompanian de ejemplos
originales aplicados a Wattpad.

\begin{enumerate}

\item \textbf{Requisito (3.1.19).}
Declaracion que expresa una necesidad
y las condiciones o restricciones
relacionadas con ella.

\textit{Ejemplo:} Wattpad debera permitir
buscar historias mediante su titulo.

\item \textbf{Ingenieria de requisitos (3.1.21).}
Disciplina encargada de identificar,
analizar, documentar, verificar,
validar y gestionar requisitos.

\textit{Ejemplo:} Analizar las necesidades
de lectores y escritores para definir
las funciones de Wattpad.

\item \textbf{Parte interesada (3.1.28).}
Persona u organizacion que posee
intereses o expectativas relacionados
con un sistema.

\textit{Ejemplo:} Los lectores y escritores
son partes interesadas de Wattpad.

\item \textbf{Usuario (3.1.35).}
Persona o grupo que interactua
con un sistema o se beneficia de el.

\textit{Ejemplo:} Un lector que consulta
capitulos desde su cuenta de Wattpad.

\item \textbf{Restriccion (3.1.7).}
Limitacion externa que condiciona
el sistema, su diseno o desarrollo.

\textit{Ejemplo:} Exigir compatibilidad
con una version minima de Android
en una propuesta de desarrollo.

\item \textbf{Condicion (3.1.6).}
Circunstancia medible que determina
cuando se aplica un requisito.

\textit{Ejemplo:} Cuando un usuario
autenticado vota por un capitulo,
el sistema registra el voto.

\item \textbf{Atributo (3.1.2).}
Propiedad de una entidad que puede
distinguirse cualitativa o
cuantitativamente.

\textit{Ejemplo:} El tiempo de respuesta
de la busqueda medido en segundos.

\item \textbf{Verificacion de requisitos (3.1.26).}
Revision que confirma que los requisitos
estan correctamente formulados.

\textit{Ejemplo:} Comprobar que el
requisito de votacion sea claro,
singular y verificable.

\item \textbf{Validacion de requisitos (3.1.25).}
Confirmacion de que los requisitos
representan las necesidades reales
de las partes interesadas.

\textit{Ejemplo:} Consultar a lectores
si las funciones de busqueda y
biblioteca satisfacen sus necesidades.

\item \textbf{Trazabilidad de requisitos (3.1.23).}
Identificacion de las relaciones
entre requisitos y sus necesidades
o elementos derivados.

\textit{Ejemplo:} Relacionar un requisito
de votacion con la necesidad de
interaccion de los lectores.

\end{enumerate}


\section{Matriz de requisitos de Wattpad}

Los siguientes requisitos constituyen una
propuesta academica basada en las
funcionalidades observadas de Wattpad.
Los umbrales de calidad son objetivos
propuestos y no mediciones reales
del funcionamiento de la plataforma.

\small
\begin{longtable}{|p{1.3cm}|p{2cm}|p{5.2cm}|p{5cm}|}
\hline
\textbf{ID} & \textbf{Tipo} &
\textbf{Requisito} &
\textbf{Criterio de aceptacion} \\ \hline
\endfirsthead

\hline
\textbf{ID} & \textbf{Tipo} &
\textbf{Requisito} &
\textbf{Criterio de aceptacion} \\ \hline
\endhead

RF-01 & Funcional &
El sistema debera mostrar historias cuyo
titulo coincida con el texto ingresado
en la busqueda. &
Se muestra un resultado correspondiente
al titulo de prueba. \\ \hline

RF-02 & Funcional &
El sistema debera mostrar el contenido
del capitulo seleccionado. &
Se visualiza el contenido del capitulo
elegido. \\ \hline

RF-03 & Funcional &
El sistema debera mostrar los comentarios
del parrafo seleccionado. &
Se visualizan los comentarios asociados
al parrafo de prueba. \\ \hline

RF-04 & Funcional &
El sistema debera registrar el voto
del usuario autenticado en un capitulo. &
El voto queda registrado y se muestra
el estado correspondiente. \\ \hline

RF-05 & Funcional &
El sistema debera mostrar las historias
guardadas en la biblioteca del usuario. &
La historia guardada aparece
en la biblioteca. \\ \hline

RI-01 & Interfaz &
La interfaz de busqueda debera mostrar
un campo para ingresar el titulo
de una historia. &
El campo acepta el texto ingresado. \\ \hline

RI-02 & Interfaz &
La interfaz de lectura debera mostrar
un control con la etiqueta Votar. &
El control aparece en la pantalla
de lectura. \\ \hline

RU-01 & Usabilidad &
Al menos 9 de 10 participantes deberan
encontrar una historia por titulo
en 60 segundos sin asistencia. &
Tasa de exito mayor o igual al 90\%
en la prueba definida. \\ \hline

RU-02 & Usabilidad &
Al menos 9 de 10 participantes deberan
guardar una historia en su biblioteca
sin asistencia. &
Tasa de exito mayor o igual al 90\%
en la prueba definida. \\ \hline

RH-01 & Factores humanos &
La interfaz debera mostrar una etiqueta
textual que identifique el estado
de voto registrado. &
Al menos 9 de 10 participantes
identifican el estado correctamente. \\ \hline

RC-01 & Calidad &
El sistema debera mostrar resultados
de busqueda en un maximo de 3 segundos
en condiciones de prueba definidas. &
Al menos 95 de 100 busquedas
cumplen el tiempo establecido. \\ \hline

RC-02 & Calidad &
El sistema debera mostrar el contenido
de un capitulo en un maximo de 3 segundos
en condiciones de prueba definidas. &
Al menos 95 de 100 aperturas
cumplen el tiempo establecido. \\ \hline

RC-03 & Calidad &
El servicio de lectura debera alcanzar
una disponibilidad mensual minima
del 99,5\%. &
Disponibilidad mensual mayor o
igual al 99,5\%. \\ \hline

RC-04 & Calidad &
El sistema debera impedir modificaciones
no autorizadas de capitulos
pertenecientes a otra cuenta. &
Se rechazan los 20 intentos
no autorizados de prueba. \\ \hline

RC-05 & Calidad &
El sistema debera mostrar una
confirmacion visual despues
de registrar un voto valido. &
Se muestra confirmacion en
20 de 20 votos validos. \\ \hline

\end{longtable}
\normalsize

\subsection{Clasificacion de calidad segun ISO/IEC 25010:2023}

\begin{center}
\small
\begin{tabular}{|p{1.4cm}|p{4.7cm}|p{5.5cm}|}
\hline
\textbf{ID} &
\textbf{Caracteristica} &
\textbf{Subcaracteristica} \\ \hline

RC-01 & Eficiencia de desempeno &
Comportamiento temporal \\ \hline

RC-02 & Eficiencia de desempeno &
Comportamiento temporal \\ \hline

RC-03 & Fiabilidad &
Disponibilidad \\ \hline

RC-04 & Seguridad &
Integridad \\ \hline

RC-05 & Capacidad de interaccion &
Autodescriptividad \\ \hline

\end{tabular}
\end{center}

Las metricas y los criterios de aceptacion
permiten planificar la verificacion de
los requisitos. No representan pruebas
ejecutadas sobre la infraestructura
interna de Wattpad.

\end{document}
